\documentclass[aps,prb,reprint,superscriptaddress,nofootinbib]{revtex4-2}

\usepackage[T1]{fontenc}
\usepackage[utf8]{inputenc}
\usepackage{amsmath,amssymb,amsthm,bm}
\usepackage{mathtools}
\usepackage{physics}
\usepackage{hyperref}

\newtheorem{definition}{Definition}
\newtheorem{proposition}{Proposition}
\newtheorem{lemma}{Lemma}
\newtheorem{remark}{Remark}

\newcommand{\Liouv}{\mathcal{L}}
\newcommand{\Hh}{\mathcal{H}}
\newcommand{\G}{\mathcal{G}}
\newcommand{\Pc}{\mathcal{P}}
\newcommand{\Qc}{\mathcal{Q}}
\newcommand{\Cc}{\mathcal{C}}
\newcommand{\Fc}{\mathcal{F}}
\newcommand{\Oc}{\mathcal{O}}
\newcommand{\dif}{\mathrm{d}}

\begin{document}

\title{Hierarchical Closure of the Equation-of-Motion Chain for the Reduced BCS Hamiltonian}

\author{A. O. C. Baeta}
\affiliation{Independent Researcher}

\date{\today}

\begin{abstract}
We formulate the equation-of-motion treatment of the reduced BCS Hamiltonian as a hierarchy of operator subspaces generated by repeated action of the Liouvillian. The purpose is not to rederive the BCS gap equation from a mean-field ansatz, but to isolate the algebraic point at which a closed anomalous propagator can be obtained from the exact Zubarev hierarchy. We define a hierarchy depth, construct the first nontrivial levels explicitly, and show that a second-order closure in the pairing interaction is obtained by repeated use of fermionic anticommutation relations together with a controlled projection onto the pair sector. The resulting closed Green-function equation has the form of a Dyson equation with a pairing self-energy. Mean-field BCS theory then appears as a rank-one projection of the hierarchy rather than as its starting assumption. The manuscript also identifies which elements are established algebraic consequences and which should be regarded as a proposed recursive formalism requiring further analysis.
\end{abstract}

\maketitle

\section{Introduction}

The reduced BCS Hamiltonian is usually introduced through a variational paired state or through a mean-field decoupling of the interaction term. This route is powerful and historically decisive, but it hides a structural question: how much of the anomalous propagator and the gap equation follows directly from the operator algebra generated by the Hamiltonian before any mean-field factorization is imposed?

Equation-of-motion (EOM) techniques, especially the double-time Green-function formalism of Zubarev, provide a natural setting for this question~\cite{Zubarev1960}. Acting with the Liouvillian on a chosen operator generates new composite operators, whose Green functions obey coupled equations. Without approximation this process produces an infinite hierarchy. Closures of such hierarchies occur throughout many-body theory, from BBGKY kinetic equations to projection-operator approaches and Hubbard-like EOM schemes~\cite{Yvon1935,BornGreen1946,Kirkwood1946,Bogoliubov1962,Mori1965,Zwanzig1960,Hubbard1963}. The present work asks whether the reduced BCS problem admits a useful closure interpretation in which the standard anomalous self-energy is recovered as the lowest nontrivial projection of a larger recursive structure.

The central claim developed here is deliberately modest. We show that the EOM chain generated by the reduced BCS Hamiltonian admits an explicit hierarchy of operator subspaces. At order \(V^2\), where \(V\) denotes the pairing interaction kernel, the chain can be closed over the Nambu pair sector by fermionic anticommutation identities and a projection onto pair-conserving contractions. This reproduces a Dyson-type anomalous propagator. We do not claim that this finite closure is exact for the full interacting hierarchy. Rather, we identify it as a controlled second-order closure and propose a language for higher recursive closures.

\section{Hamiltonian and Green Functions}

We consider the reduced BCS Hamiltonian
\begin{equation}
H = H_0 + H_{\mathrm{int}},
\end{equation}
with
\begin{equation}
H_0 = \sum_{k\sigma}\xi_k c^\dagger_{k\sigma}c_{k\sigma},
\end{equation}
and
\begin{equation}
H_{\mathrm{int}}
= \sum_{kk'} V_{kk'}
c^\dagger_{k\uparrow}c^\dagger_{-k\downarrow}
c_{-k'\downarrow}c_{k'\uparrow}.
\label{eq:bcs-hint}
\end{equation}
The sign convention is absorbed into \(V_{kk'}\); attractive pairing corresponds to the usual negative pairing channel.

For fermionic operators \(A,B\), the retarded double-time Green function is
\begin{equation}
G_{AB}(\omega)
= \langle\langle A;B\rangle\rangle_\omega
= -i\int_0^\infty \dif t\, e^{i(\omega+i0^+)t}
\langle \{A(t),B(0)\}\rangle .
\end{equation}
The Zubarev EOM is
\begin{equation}
\omega \langle\langle A;B\rangle\rangle_\omega
= \langle\{A,B\}\rangle
+ \langle\langle [A,H];B\rangle\rangle_\omega .
\label{eq:zubarev}
\end{equation}
It is useful to introduce the Liouvillian
\begin{equation}
\Liouv A = [A,H],\qquad
\Liouv_0 A=[A,H_0],\qquad
\Liouv_I A=[A,H_{\mathrm{int}}].
\end{equation}

\begin{definition}[EOM hierarchy]
Given a seed operator \(A_0\), define the hierarchy generated by \(H\) as the sequence of operator subspaces
\begin{equation}
\Hh_n(A_0)=\mathrm{span}\{\Liouv^m A_0:0\le m\le n\}.
\end{equation}
The strict level \(n\) is the quotient contribution
\(\Hh_n/\Hh_{n-1}\), after reducing operators by canonical anticommutation relations.
\end{definition}

This definition is purely algebraic. It does not require an ansatz for the expectation values and therefore separates the growth of the operator hierarchy from later approximation.

\section{The Operator Hierarchy}

Take as seed the normal electron operator
\begin{equation}
A_0=c_{k\uparrow}.
\end{equation}
The free Liouvillian gives
\begin{equation}
\Liouv_0 c_{k\uparrow}=\xi_k c_{k\uparrow}.
\end{equation}
The interaction part produces
\begin{equation}
\Liouv_I c_{k\uparrow}
= \sum_q V_{kq} c^\dagger_{-k\downarrow}
c_{-q\downarrow}c_{q\uparrow},
\label{eq:first-level}
\end{equation}
up to the sign convention fixed in Eq.~\eqref{eq:bcs-hint}. Thus the first nontrivial level contains a three-fermion operator. Similarly, choosing the anomalous seed \(c^\dagger_{-k\downarrow}\) generates
\begin{equation}
\Liouv_I c^\dagger_{-k\downarrow}
= -\sum_q V_{qk} c^\dagger_{q\uparrow}
c^\dagger_{-q\downarrow}c_{k\uparrow}.
\end{equation}
The EOM for the Nambu spinor therefore immediately couples one-particle propagators to composite three-operator propagators.

\begin{proposition}[Hierarchy growth]
For the reduced BCS interaction, \(\Liouv_I\) maps a monomial with \(r\) fermionic factors into a finite sum of monomials with \(r+2\), \(r\), and \(r-2\) factors after canonical anticommutation reduction.
\end{proposition}

\begin{proof}
The interaction term contains four fermionic factors. Commuting a monomial through it generates products with the original \(r\) factors and four additional factors. Each contraction produced by \(\{c_i,c^\dagger_j\}=\delta_{ij}\) removes two factors. Since at least one contraction is required in the commutator with an odd fermionic monomial of the same algebra, the possible reduced degrees differ from \(r\) by even integers. For the BCS quartic vertex, the relevant reduced sectors are \(r+2\), \(r\), and \(r-2\), subject to momentum and spin constraints.
\end{proof}

The proposition is the algebraic origin of the hierarchy. In the BCS case, the \(r=1\) seed generates \(r=3\) operators, whose EOM generate one-, three-, and five-operator sectors. A closed one-particle equation is therefore not exact unless the higher sectors are eliminated or represented recursively.

\section{Perturbative Structure}

Let \(V\) denote the scale of the interaction kernel. The formal resolvent associated with Eq.~\eqref{eq:zubarev} is
\begin{equation}
G_{AB}(\omega)
= \langle\{(\omega-\Liouv)^{-1}A,B\}\rangle,
\end{equation}
where the notation suppresses the thermal inner product structure. Expanding
\begin{equation}
(\omega-\Liouv)^{-1}
= (\omega-\Liouv_0)^{-1}
+(\omega-\Liouv_0)^{-1}\Liouv_I(\omega-\Liouv_0)^{-1}
+\cdots
\end{equation}
shows that a contribution of order \(V^n\) corresponds to \(n\) insertions of \(\Liouv_I\). At order \(V^2\), the hierarchy can leave the one-particle sector and return to it through the pair channel.

\begin{definition}[Pair projection]
Let \(\Pc_{\mathrm{pair}}\) denote the algebraic projection that retains contractions returning a composite operator to the Nambu pair sector
\begin{equation}
\mathrm{span}\{c_{k\uparrow},c^\dagger_{-k\downarrow}\},
\end{equation}
and discards residual irreducible operators with three or more uncontracted fermionic factors.
\end{definition}

This projection is not a statement of exact equality in the full algebra. It defines a controlled closure prescription. Its validity must be assessed by the order retained and by comparison with omitted irreducible sectors.

\section{Closure at Order \(V^2\)}

Define the Nambu Green function
\begin{equation}
\hat{G}_k(\omega)
=
\begin{pmatrix}
\langle\langle c_{k\uparrow};c^\dagger_{k\uparrow}\rangle\rangle_\omega&
\langle\langle c_{k\uparrow};c_{-k\downarrow}\rangle\rangle_\omega\\
\langle\langle c^\dagger_{-k\downarrow};c^\dagger_{k\uparrow}\rangle\rangle_\omega&
\langle\langle c^\dagger_{-k\downarrow};c_{-k\downarrow}\rangle\rangle_\omega
\end{pmatrix}.
\end{equation}
The exact first EOM has the schematic form
\begin{equation}
(\omega-\xi_k)G_k(\omega)
=1+\sum_q V_{kq}
\langle\langle c^\dagger_{-k\downarrow}c_{-q\downarrow}c_{q\uparrow};
c^\dagger_{k\uparrow}\rangle\rangle_\omega .
\label{eq:normal-eom-composite}
\end{equation}
The anomalous equation has an analogous composite propagator. Acting once more with the Liouvillian on the composite operator produces terms that contain contractions proportional to pair expectation values
\begin{equation}
\Phi_q=\langle c_{-q\downarrow}c_{q\uparrow}\rangle,
\end{equation}
and residual higher operators. Applying \(\Pc_{\mathrm{pair}}\) gives
\begin{equation}
\Pc_{\mathrm{pair}}\Liouv_I
\left(c^\dagger_{-k\downarrow}c_{-q\downarrow}c_{q\uparrow}\right)
\sim
\Delta_k c^\dagger_{-k\downarrow}
+\text{normal energy-shift terms},
\end{equation}
where
\begin{equation}
\Delta_k=-\sum_q V_{kq}\Phi_q.
\label{eq:gap-definition}
\end{equation}

\begin{lemma}[Second-order pair return]
Under \(\Pc_{\mathrm{pair}}\), the sequence
\begin{equation}
c_{k\uparrow}
\xrightarrow{\Liouv_I}
c^\dagger_{-k\downarrow}c_{-q\downarrow}c_{q\uparrow}
\xrightarrow{\Liouv_I}
c^\dagger_{-k\downarrow}
\end{equation}
is nonzero only through contractions in the pair channel and is of order \(V^2\) in the EOM hierarchy.
\end{lemma}

\begin{proof}
The first arrow is Eq.~\eqref{eq:first-level} and carries one factor of \(V\). The second action of \(\Liouv_I\) produces quartic interaction operators multiplied by the three-fermion monomial. Canonical anticommutation reduction contains contractions of the annihilation pair \(c_{-q\downarrow}c_{q\uparrow}\) with the corresponding creation pair in \(H_{\mathrm{int}}\), leaving a single operator \(c^\dagger_{-k\downarrow}\), weighted by a pair amplitude. All terms with uncontracted three- or five-operator content are outside \(\Pc_{\mathrm{pair}}\). Hence the projected return is second order in the interaction insertions and lands in the anomalous Nambu sector.
\end{proof}

The closed projected EOM becomes
\begin{equation}
\begin{pmatrix}
\omega-\xi_k & -\Delta_k\\
-\Delta_k^\ast & \omega+\xi_k
\end{pmatrix}
\hat{G}_k(\omega)=\hat{1},
\end{equation}
with solution
\begin{equation}
\hat{G}_k(\omega)
=
\frac{1}{\omega^2-E_k^2}
\begin{pmatrix}
\omega+\xi_k & \Delta_k\\
\Delta_k^\ast & \omega-\xi_k
\end{pmatrix},
\qquad
E_k=\sqrt{\xi_k^2+|\Delta_k|^2}.
\end{equation}
Using the anomalous component gives the self-consistency relation
\begin{equation}
\Delta_k
=-\sum_q V_{kq}
\frac{\Delta_q}{2E_q}\tanh\left(\frac{\beta E_q}{2}\right),
\end{equation}
which is the BCS gap equation. In the present formulation this equation is the fixed point of the second-order projected hierarchy, not an initial mean-field assumption.

\section{Dyson Form}

Let
\begin{equation}
\hat{G}_{0,k}^{-1}(\omega)=
\begin{pmatrix}
\omega-\xi_k&0\\
0&\omega+\xi_k
\end{pmatrix}.
\end{equation}
The projected closure can be written
\begin{equation}
\hat{G}_k^{-1}(\omega)
=\hat{G}_{0,k}^{-1}(\omega)-\hat{\Sigma}_k(\omega),
\end{equation}
where, at the lowest pair closure level,
\begin{equation}
\hat{\Sigma}_k=
\begin{pmatrix}
0&\Delta_k\\
\Delta_k^\ast&0
\end{pmatrix}.
\end{equation}
This is a Dyson equation, but the route to it has been algebraic: the self-energy is the representation of eliminated hierarchy levels under the closure map.

\begin{definition}[Closure map]
A closure map of order \(n\) is a linear or nonlinear map
\begin{equation}
\Cc_n:\Hh_{n+1}\to \Hh_n
\end{equation}
that replaces the first omitted hierarchy level by operators in retained levels, with an explicitly stated error sector \(\Qc_n=1-\Pc_n\).
\end{definition}

\begin{proposition}[Dyson representation of a projected hierarchy]
If the EOM hierarchy is split into retained and omitted sectors by \(\Pc+\Qc=1\), and if the omitted sector is represented by a closure map \(\Cc\), then the retained Green function satisfies a Dyson-type equation with self-energy equal to the induced action of \(\Cc\) on the retained resolvent.
\end{proposition}

\begin{proof}
Write the Liouvillian in block form with respect to \(\Pc\oplus\Qc\). Eliminating the \(\Qc\) block by a Schur complement gives
\begin{equation}
G_P^{-1}(\omega)
=\omega-\Liouv_{PP}
-\Liouv_{PQ}(\omega-\Liouv_{QQ})^{-1}\Liouv_{QP}.
\end{equation}
Replacing the exact \((\omega-\Liouv_{QQ})^{-1}\) by the closure prescription defines the projected self-energy. This is the algebraic form of Dyson's equation.
\end{proof}

\section{Mean Field as Projection}

The usual BCS mean-field Hamiltonian is
\begin{equation}
H_{\mathrm{MF}}=
\sum_{k\sigma}\xi_k c^\dagger_{k\sigma}c_{k\sigma}
-\sum_k \left(
\Delta_k c^\dagger_{k\uparrow}c^\dagger_{-k\downarrow}
+\Delta_k^\ast c_{-k\downarrow}c_{k\uparrow}
\right).
\end{equation}
In the present formalism, \(H_{\mathrm{MF}}\) is the effective Hamiltonian whose Liouvillian reproduces the projected action of the full Liouvillian on the Nambu sector:
\begin{equation}
\Pc_{\mathrm{pair}}\Liouv \Pc_{\mathrm{pair}}
\quad\longleftrightarrow\quad
\Liouv_{\mathrm{MF}}.
\end{equation}
Thus mean field is a projection of the hierarchy. This perspective is useful because it makes the approximation explicit: the omitted content is not hidden in a factorization step, but assigned to \(\Qc_{\mathrm{pair}}\), the residual hierarchy.

\section{Recursive Algebra and Open Formalization}

The preceding sections establish a finite second-order closure. A stronger program would define a recursive sequence
\begin{equation}
\Hh_0\subset\Hh_1\subset\cdots,\qquad
\Cc_n:\Hh_{n+1}\to\Hh_n,
\end{equation}
and study convergence or stability of the induced self-energies
\begin{equation}
\Sigma^{(n)}(\omega)
=\Liouv_{0,1}(\omega-\Liouv_{1,1}^{(n)})^{-1}\Liouv_{1,0}.
\end{equation}
At this stage this is a proposed formalization, not a theorem. Possible mathematical settings include Banach-space closures of resolvents, continued-fraction representations of Liouvillian dynamics, or homological descriptions of contractions and residual operator sectors. These directions require additional proof of domain control, norm estimates, and uniqueness of the closure map.

\begin{definition}[Hierarchical information content]
For a chosen projection \(\Pc_n\), define the residual hierarchy weight
\begin{equation}
I_n(\omega)=
\norm{\Qc_n(\omega-\Liouv)^{-1}\Pc_n}^{2},
\end{equation}
in any specified operator norm for which the resolvent is well-defined.
\end{definition}

\begin{remark}
The quantity \(I_n\) is introduced only as a possible diagnostic of closure quality. No monotonicity, universality, or variational principle is claimed here.
\end{remark}

\section{Relation to Earlier Work}

BCS theory introduced the microscopic paired ground state and the self-consistent gap equation~\cite{Bardeen1957}. Gor'kov reformulated superconductivity in terms of Green functions and anomalous propagators~\cite{Gorkov1958}. Zubarev's double-time Green functions provided the EOM framework used here~\cite{Zubarev1960}. Dyson's equation supplies the general language in which eliminated degrees of freedom appear as self-energies~\cite{Dyson1949}. Fetter and Walecka and Mahan give standard many-body accounts of these ideas~\cite{FetterWalecka1971,Mahan2000}. Mori-Zwanzig projection methods and the BBGKY hierarchy show, in different settings, how exact dynamical equations generate higher-order structures requiring closure~\cite{Mori1965,Zwanzig1960,Yvon1935,BornGreen1946,Kirkwood1946,Bogoliubov1962}. Hubbard's EOM treatment of correlated electrons is another important precedent for operator hierarchies closed by physically motivated projections~\cite{Hubbard1963}.

The intended novelty of the present manuscript is therefore not the existence of the BCS gap equation, the anomalous Green function, or the Dyson equation. The proposed contribution is the explicit organization of the reduced BCS EOM as a hierarchy in which the BCS mean-field equations arise as the lowest nontrivial pair-sector closure.

\section{Conclusions}

We have recast the EOM derivation of BCS superconductivity as a hierarchical closure problem. The exact Liouvillian hierarchy generated by the reduced BCS Hamiltonian grows beyond the Nambu sector after one interaction insertion. A second interaction insertion allows a return to the pair sector through fermionic contractions. Projecting this return yields the standard anomalous self-energy and the BCS gap equation. This makes mean-field BCS theory appear as a controlled pair-sector projection of the EOM hierarchy rather than as the starting point of the analysis.

The next mathematical task is to define higher closure maps with precise error estimates. The next physical task is to determine whether higher residual sectors encode corrections that can be organized systematically beyond BCS mean field.

\begin{acknowledgments}
This manuscript was prepared as a working draft for a hierarchy-based reformulation of the BCS equation-of-motion derivation.
\end{acknowledgments}

\appendix

\section{Explicit Fermionic Commutators}

This appendix records the algebra behind the hierarchy statements used in the main text. Let
\begin{equation}
P_k=c_{-k\downarrow}c_{k\uparrow},\qquad
P_k^\dagger=c_{k\uparrow}^\dagger c_{-k\downarrow}^\dagger .
\end{equation}
The interaction is written as
\begin{equation}
H_{\mathrm{int}}=\sum_{pq}V_{pq}P_p^\dagger P_q .
\end{equation}
Using the canonical anticommutation relations,
\begin{equation}
\{c_{i},c_{j}^\dagger\}=\delta_{ij},\qquad
\{c_i,c_j\}=\{c_i^\dagger,c_j^\dagger\}=0,
\end{equation}
one obtains
\begin{align}
[c_{k\uparrow},P_p^\dagger]
&=[c_{k\uparrow},c_{p\uparrow}^\dagger c_{-p\downarrow}^\dagger]\nonumber\\
&=\delta_{kp}c_{-p\downarrow}^\dagger ,
\end{align}
and
\begin{equation}
[c_{k\uparrow},P_q]=0 .
\end{equation}
Therefore
\begin{align}
[c_{k\uparrow},H_{\mathrm{int}}]
&=\sum_{pq}V_{pq}[c_{k\uparrow},P_p^\dagger]P_q\nonumber\\
&=\sum_qV_{kq}c_{-k\downarrow}^\dagger c_{-q\downarrow}c_{q\uparrow}.
\end{align}
This is the first explicit step from the one-fermion sector to the three-fermion sector. Similarly,
\begin{align}
[c_{-k\downarrow}^\dagger,P_q]
&=[c_{-k\downarrow}^\dagger,c_{-q\downarrow}c_{q\uparrow}]\nonumber\\
&=-\delta_{kq}c_{q\uparrow},
\end{align}
which gives
\begin{equation}
[c_{-k\downarrow}^\dagger,H_{\mathrm{int}}]
=-\sum_pV_{pk}P_p^\dagger c_{k\uparrow}.
\end{equation}
The signs depend only on the ordering convention chosen for \(P_k\); physical observables are unaffected once the convention is used consistently.

\section{Second-Order Pair-Sector Return}

Consider the composite operator generated after one interaction insertion,
\begin{equation}
X_{kq}=c_{-k\downarrow}^\dagger P_q
=c_{-k\downarrow}^\dagger c_{-q\downarrow}c_{q\uparrow}.
\end{equation}
The next interaction commutator is
\begin{equation}
[X_{kq},H_{\mathrm{int}}]
=\sum_{pl}V_{pl}[X_{kq},P_p^\dagger P_l].
\end{equation}
Using
\begin{equation}
[X,AB]=[X,A]B+A[X,B],
\end{equation}
one separates terms that either increase the operator degree or contract back into the Nambu sector. The pair-return contribution is obtained from the contraction of \(P_q\) with \(P_p^\dagger\):
\begin{align}
P_qP_p^\dagger
&=c_{-q\downarrow}c_{q\uparrow}
c_{p\uparrow}^\dagger c_{-p\downarrow}^\dagger\nonumber\\
&=\delta_{qp}(1-n_{q\uparrow}-n_{-q\downarrow})
+\text{normal-ordered four-fermion terms}.
\end{align}
Substitution into \(X_{kq}P_p^\dagger P_l\) gives, after retaining the pair-sector contraction,
\begin{equation}
\Pc_{\mathrm{pair}}[X_{kq},H_{\mathrm{int}}]
\propto
\sum_lV_{ql}
\langle P_l\rangle c_{-k\downarrow}^\dagger
+\text{normal shifts}.
\end{equation}
Summing over the first interaction vertex gives
\begin{equation}
\sum_qV_{kq}\Pc_{\mathrm{pair}}[X_{kq},H_{\mathrm{int}}]
\propto
-\Delta_k c_{-k\downarrow}^\dagger,
\end{equation}
with
\begin{equation}
\Delta_k=-\sum_qV_{kq}\langle P_q\rangle.
\end{equation}
Thus the anomalous Nambu operator is reached only after two interaction insertions:
\begin{equation}
c_{k\uparrow}
\xrightarrow{\Liouv_I}
X_{kq}
\xrightarrow{\Pc_{\mathrm{pair}}\Liouv_I}
c_{-k\downarrow}^\dagger .
\end{equation}
The omitted part consists of residual normal-ordered three- and five-fermion operators. These residual operators are precisely the first error sector of the lowest closure.

\section{Finite-Size Pair Diagnostics}

For a finite Hamiltonian with exact particle-number conservation, the anomalous expectation value
\begin{equation}
\langle P_j\rangle
\end{equation}
vanishes unless a symmetry-breaking source is introduced. Exact diagonalization should therefore not be compared directly with the broken-symmetry BCS amplitude. The number-conserving diagnostic used in the computational companion is the pair-density matrix
\begin{equation}
C_{jl}=\langle P_j^\dagger P_l\rangle .
\end{equation}
The lowest pair-sector closure predicts a rank-one approximation
\begin{equation}
C_{jl}^{\mathrm{BCS}}\propto \phi_j^\ast\phi_l,\qquad
\phi_j=\frac{\Delta}{2\sqrt{\xi_j^2+|\Delta|^2}},
\end{equation}
up to an optimal scalar normalization. The reported residual is
\begin{equation}
R_C=
\frac{\|C-C^{\mathrm{BCS}}\|_F}{\|C\|_F}.
\end{equation}
This residual is not a theorem of convergence. It is a numerical diagnostic of how much pair-sector information is missed by the rank-one closure in a finite system.

\bibliographystyle{apsrev4-2}
\bibliography{references}

\end{document}
